\paragraph{Three exchanges and the removal of the twist}
Apply the exchange of Section~\ref{sec:exchange-proof} to kernels whose real
data are products of determinants. The real vectors obtained by exchange
can then be replaced with explicit polynomials. Split the positive coordinates
as \(\mathbb C^p\oplus\mathbb C^{d_+}\), with standard frame \((I_p\ \ 0)\).
For \(z\in U(1)\) act by \(z^{-2}\) on the first \(p\) coordinates,
by \(z^{-1}\) on the next \(d_+\), and by \(z^{-1}I_{d_+}\) on the partner.
With the determinant factors in the Weil action included, all variables
outside the last square block acquire positive degree in \(z\). The circle
average is consequently evaluation with those variables set to zero.
A determinant-weighted average over \(U(d_+)\) selects the component with
\(P(kg)=\det(k)^aP(g)\), spanned by \(\det^a\).
The identity holds first for unitary translates of the standard frame.
Polarization gives it for products with arbitrary auxiliary rows.
This allows the exchange to use the determinants of stacked rows and the
explicit matched polynomial on the column side. Density extends the
averaging argument from polynomials to Fock vectors under the continuous
functionals in the two integrals.
\leanlinks{\leansrc{New/Determinant/Exchange/ExplicitMatching}{Circle averaging, projection onto determinant powers, and polarization}}

Set \(c=\chi_M^j\vartheta\), a character of weight zero. The fixed
character \(\chi_M\) has residue-field order greater than \(8m\) at the
chosen places. The first exchange uses a definite space \(J\) of rank
\(m\) and the zero space as partners, with finite data \(\chi\circ\det\)
and \(1\). Integration on the zero space gives a constant, so a scalar
makes the square-kernel class sum on \(J\) equal to \(1\). Its squared
modulus is proportional to \(L_*(m/2;c\upsilon^{m/2})^{-1}\).
Here \(L_*(u;\nu')\) denotes \eqref{eq:exchange-lstar} with character \(\nu'\).
Conjugate this identity for use as a kernel from \(J^{\mathrm{opp}}\).
In the second exchange, let \(S_0\) have signature \((2m,m)\).
The central space is \(J\perp S_0\), and its partners are the definite
space of rank \(3m\) and \(J^{\mathrm{opp}}\).
The kernel on the large side splits into the kernels for \(J\) and \(S_0\),
so the preceding class sum removes the \(J\) factor.
The remaining form \(\phi^\circ\) is the lift of \(1\) from \(J^{\mathrm{opp}}\)
to \(S_0\), independent of the twist. The second exchange contributes
an additional \(L_*(m;c^{-1}\upsilon^m)^{-1}\) to the squared modulus.
The determinant character from the change of finite variables cancels
that in the Weil action, giving standard finite level for \(\phi^\circ\).
The real reproducing identity gives holomorphy and rapid decrease.
To see nonvanishing, let \(v\) be the Fock vector given by the
\(2m\)-th power of the determinant in the negative block at each place.
On the compact real group of \(J^{\mathrm{opp}}\), the determinant from the
block action cancels the factor in the Weil action. Thus
\(\langle v,\omega(u,1)v\rangle=\langle v,v\rangle\).
The reproducing identity with the finite vacuum would make its average
zero if the lift vanished. Since \(v\ne0\) and the volume is positive,
this is impossible.
\leanlinks{\leansrc{New/OneLift/ComparisonOperator/ThreeExchanges}{The square-kernel class sum and the twist-independent lift}}

For the third exchange use the central space obtained by adjoining
the standard positive \(E^{3m}\) to
\[
 W'=W_j\perp(E^j)^{\mathrm{opp}}\perp E^{3m-j}\perp E^{6m+1}.
\]
The large partner \(V''\) is definite of the same rank \(r+9m+1\)
as \(W'\), and the small partner is \(S_0\). The half-difference of
ranks is \(z_*=(r+6m+1)/2=7m/2+k/2\).
For the finite datum use the trace character specified by a matrix
\(\Lambda\). Average it over \(\mathrm{GL}_{3m}\) modulo \(\pi_M^2\)
to make it invariant under right changes of basis in \(E^{3m}\).
Sum this exchange over the classes of \(E^{3m}\) and insert the result
of the first two exchanges. Dividing by the order of the finite group
leaves the integral \(F_i(g)\) of \(\phi^\circ\) against the remaining
kernel. Here \(i\) records \(\Lambda\), an auxiliary class, and the fixed
rows. On the row side, integrating the square block turns the double
class sum into the theta series \(\Xi_{\vartheta,i}\) on \(W_j\)
with column data. Thus
\begin{equation}\label{eq:comparison-period-factor}
 c'_\vartheta\Xi_{\vartheta,i}=C_i F_i,\qquad
 |c'_\vartheta|^2=C'/\mathcal H_j(\vartheta).
\end{equation}
Each \(C_i\) is a nonzero normalization independent of the twist.
The three factors of \(\mathcal H_j\) are \eqref{eq:exchange-lstar} at
\(m/2,m,z_*\), with the inverses of
\(\nu_-=\upsilon^{-m/2}\chi_M^{-j}\),
\(\nu_+=\upsilon^{-m}\chi_M^j\), and
\(\nu_*=\mu_j^0\upsilon^{j+(t-11m-1)/2}\) after multiplication by
\(\vartheta^{-1},\vartheta,\vartheta\), respectively.
\leanlinks{\leansrc{New/OneLift/ComparisonOperator/ThreeExchanges}{The third exchange and the three critical pairs}
  \leanlinksep \leansrc{New/OneLift/ComparisonOperator/FirstFactors}{Normalization of the first block's class sum}}

\paragraph{Gluing and the period map}
Use \(V_j=V''\perp V^\circ\). The glued lattice is generated by pairs of
standard inclusion columns, the sublattice in \(V''\) on which the rows
specified by \(\Lambda\) vanish modulo \(\pi_M\), and \(\pi_M\) times the
lattice of \(V^\circ\). For a basis matrix \(A\), \(\pi_M A^{-1}\) is
integral. Rewriting the congruence on lift columns in this lattice
separates it into finitely many conditions modulo \(\pi_M^2\).
For each condition the Weil kernel factors, and the adjoint of the
Cartan product separates the pairing of values. Using
\(\Xi_{\vartheta,i}\) and the kernel \(\theta^\circ\) for \(V^\circ\),
take the two class sums first. For \(f\in F_j\) and a tuple \(z=(z_i)\)
one obtains
\[
 c'_\vartheta\langle z,G_{j,\vartheta}^*u_{j,\vartheta} f\rangle
    =\sum_i d_i\int_{[U(W_j)]}
      f(g)F_i(g)\langle z_i,\theta^\circ(g,\cdot)\rangle\,dg .
\]
The nonzero \(d_i\) are the fixed normalization factors from the
level indices and \eqref{eq:comparison-period-factor}. Let \(T\) be the
map defined by the right side.
\leanlinks{\leansrc{New/OneLift/ComparisonOperator/GluedForms}{The glued lattice and the product decomposition of the kernel}
  \leanlinksep \leansrc{New/OneLift/ComparisonOperator/GluingMaps}{The inner-product formula for the period map}}

Injectivity of \(T\) follows from a kernel approximation. Embed \(W_j\)
in the space with padded signs and use the product \(P\) of minors in
matrix variables of the same sign. Its Fock matrix coefficient \(c_P\)
has \(c_P(1,1)>0\), and unitarity and Cauchy--Schwarz bound it everywhere.
Take a peak \(g_1\) of \(\phi^\circ\), where \(\phi^\circ(g_1)\ne0\).
With the determinant weight \(a(h')\) on the real group set
\[
 F_0(h)=\int_{U(S_0)(E_\infty)}
            \phi^\circ(g_1h')\,c_P(h',h)\,a(h')\,dh'.
\]
Holomorphic reproduction gives \(F_0(1)=J_0\phi^\circ(g_1)c_P(1,1)\),
where \(J_0>0\) is the integral of the squared modulus of the relevant
real minors. As in Section~\ref{sec:spanning}, finite data isolate the
term from the specified rational columns, and a lattice sum bounds
the other terms. For the polynomial of degree at most \(b^{3/2}\),
the approximation applies at points \(g\) with first minimum at least
\((5b^{24})^{-1/2}\). Put \(n_*=2ft^3\). The error in the kernel is at most
\[
 b^{-20n_*}c_P(1,1)(\|h'\|_{\mathrm{op}}\|h\|_{\mathrm{op}})^{6n_*}.
\]
Integrating against \(\phi^\circ(g_1h')a(h')\) contributes a factor
at most \(2b^{n_*}J_0|\phi^\circ(g_1)|\), by the integral of the real minors.
Hence some linear combination \(F=\sum_i z_i' F_i\) satisfies
\[
 |F(gh)-F_0(h)|\le 2b^{-19n_*}|F_0(1)|\,\|h\|_{\mathrm{op}}^{6n_*}.
\]
If \(Tf=0\), then all the lifts to \(V^\circ\) of \(fF_i\) vanish.
Finite level and the weight gaps ensure that, if \(f\ne0\), the first
minimum at a peak of \(\bar f\) is at least \((5b^{24})^{-1/2}\).
Applying the reproducing formula to the vanishing lifts at such a peak,
and inserting the error bound, gives
\[
 |F_0(1)|\le (4b^{n_*}+1)\,2b^{-19n_*}|F_0(1)|.
\]
The coefficient on the right is less than \(1\) since \(b^{n_*}>10\).
As \(F_0(1)\ne0\), this proves that \(T\) is injective.
\leanlinks{\leansrc{New/OneLift/ComparisonOperator/ProductsSpan}{Kernel approximation and reproduction separate the periods}}

We now have \(G_{j,\vartheta}^*u_{j,\vartheta}=(c'_\vartheta)^{-1}T\).
Let \(a_j\) be the product of the finite level group order and the square
of the hook number appearing in the comparison operator.
Then \(U_j^\circ=a_j|c'_1|^{-2}T^*T\) is positive definite.
Equation~\eqref{eq:comparison-period-factor} gives
\[
 u_{j,\vartheta}^*U_{j,\vartheta}u_{j,\vartheta}
   =a_j\,|c'_\vartheta|^{-2}T^*T
   =\frac{\mathcal H_j(\vartheta)}{\mathcal H_j(1)}\,U_j^\circ .
\]
For a false lift, its descent vanishes, and the ramified character kills
the hyperplane sums from nonprimitive columns. Pulling the form back by
\(G_{j,\vartheta}^*\) therefore makes every pairing with the \(V^\circ\)
kernel in the period formula vanish. Thus \(U_{j,\vartheta}u_{j,\vartheta} f\)
is orthogonal to the false lifts.
\leanlinks{\leansrc{New/OneLift/ComparisonOperator/GluingMaps}{Orthogonality to false lifts and injectivity of the period map}}
